\documentclass[11pt,a4paper]{article}
\usepackage[utf8]{inputenc}
\usepackage{amsmath,amsfonts,amssymb}
\usepackage{graphicx}
\usepackage{geometry}
\geometry{margin=1in}
\usepackage{booktabs}
\usepackage{hyperref}
\usepackage{float}
\usepackage{xcolor}

\definecolor{darkbg}{HTML}{0d121f}
\definecolor{lighttext}{HTML}{e2e8f0}
\definecolor{primary}{HTML}{8b5cf6}
\definecolor{secondary}{HTML}{0dd5c5}

\pagecolor{darkbg}
\color{lighttext}

\hypersetup{
    colorlinks=true,
    linkcolor=secondary,
    urlcolor=secondary
}

\title{\textbf{EE3621: Digital Signal Processing} \\ \Large Lecture 12: IIR Filter Design — Analog Prototype \& Bilinear Transform \\ \large (III B.Tech EEE)}
\author{Lecture Notes (40-Minute Plan)}
\date{}

\begin{document}

\maketitle

\section{Lecture Plan \& Objectives (40 Minutes Breakdown)}
\begin{itemize}
    \item \textbf{00:00 -- 05:00 (5 mins):} \textbf{IIR Design Philosophy} --- Borrowing from continuous-time analog filter theory.
    \item \textbf{05:00 -- 15:00 (10 mins):} \textbf{Analog Prototypes} --- Butterworth, Chebyshev (I \& II), and Elliptic filters.
    \item \textbf{15:00 -- 20:00 (5 mins):} \textbf{Impulse Invariance Method} --- Time-domain mapping and aliasing.
    \item \textbf{20:00 -- 28:00 (8 mins):} \textbf{The Bilinear Transform} --- Frequency-domain mapping and warping.
    \item \textbf{28:00 -- 35:00 (7 mins):} \textbf{Bilinear Design Procedure} --- Worked example for a 3rd-order Butterworth.
    \item \textbf{35:00 -- 40:00 (5 mins):} \textbf{Checkpoints} --- Review questions.
\end{itemize}

\noindent\rule{\textwidth}{0.5pt}

\section{IIR Design Philosophy}
Infinite Impulse Response (IIR) digital filters are typically designed by first designing a continuous-time analog prototype. The standard approach involves:
\begin{enumerate}
    \item \textbf{Pre-warping:} Convert digital frequency specs to analog frequency specs.
    \item \textbf{Analog Design:} Compute the analog transfer function $H_a(s)$.
    \item \textbf{Transformation:} Map $H_a(s)$ to the digital transfer function $H(z)$.
\end{enumerate}

\section{Butterworth Analog Prototype}
The Butterworth filter is maximally flat in the passband. 

\subsection{Magnitude Response}
The squared magnitude response is:
\begin{equation}
|H_a(j\Omega)|^2 = \frac{1}{1 + \left(\frac{\Omega}{\Omega_c}\right)^{2N}}
\end{equation}

\subsection{Derivation of Pole Locations}
Substituting $s = j\Omega$ ($\Omega = s/j$):
\begin{equation}
H_a(s) H_a(-s) = \frac{1}{1 + \left(\frac{s/j}{\Omega_c}\right)^{2N}} = \frac{1}{1 + (-1)^N \left(\frac{s}{\Omega_c}\right)^{2N}}
\end{equation}
Setting the denominator to zero:
\begin{equation}
1 + (-1)^N \left(\frac{s}{\Omega_c}\right)^{2N} = 0 \implies \left(\frac{s}{\Omega_c}\right)^{2N} = (-1)^{N-1}
\end{equation}
The stable poles (LHP) are:
\begin{equation}
s_k = \Omega_c e^{j\frac{\pi}{2N}(2k + N - 1)}, \quad \text{for } k = 1, 2, \dots, N
\end{equation}
These poles lie precisely on a circle of radius $\Omega_c$.

\subsection{Derivation of Order Formula}
Given passband attenuation $A_p$ at $\Omega_p$ and stopband attenuation $A_s$ at $\Omega_s$:
\begin{equation}
\left(\frac{\Omega_p}{\Omega_c}\right)^{2N} = A_p^2 - 1 \quad \text{and} \quad \left(\frac{\Omega_s}{\Omega_c}\right)^{2N} = A_s^2 - 1
\end{equation}
Dividing and solving for $N$:
\begin{equation}
N \geq \frac{\log_{10}\left( \frac{A_s^2 - 1}{A_p^2 - 1} \right)}{2 \log_{10}\left( \frac{\Omega_s}{\Omega_p} \right)}
\end{equation}

\section{Chebyshev and Elliptic Prototypes}
\begin{itemize}
    \item \textbf{Chebyshev Type I:} Equiripple passband, monotonic stopband. $|H(j\Omega)|^2 = \frac{1}{1 + \epsilon^2 C_N^2(\Omega/\Omega_p)}$. Poles lie on an ellipse.
    \item \textbf{Chebyshev Type II:} Monotonic passband, equiripple stopband.
    \item \textbf{Elliptic (Cauer):} Equiripple in both bands, providing the minimum possible order. Uses Jacobian elliptic functions.
\end{itemize}

\begin{table}[H]
\centering
\begin{tabular}{lllll}
\toprule
\textbf{Filter Type} & \textbf{Passband} & \textbf{Stopband} & \textbf{Phase Linearity} & \textbf{Order Required} \\
\midrule
Butterworth & Maximally Flat & Monotonic & Best & Highest \\
Chebyshev I & Equiripple & Monotonic & Moderate & Lower \\
Chebyshev II & Monotonic & Equiripple & Moderate & Lower \\
Elliptic & Equiripple & Equiripple & Non-linear & Lowest \\
\bottomrule
\end{tabular}
\caption{Comparison of Analog Prototype Filters}
\end{table}

\section{The Impulse Invariance Method}
The impulse invariance method defines $h[n] = T_s \, h_a(nT_s)$. 
An analog pole at $s_k$ maps directly to a digital pole at $z_k = e^{s_k T_s}$:
\begin{equation}
H_a(s) = \sum_{k=1}^N \frac{A_k}{s - s_k} \implies H(z) = \sum_{k=1}^N \frac{T_s A_k}{1 - e^{s_k T_s} z^{-1}}
\end{equation}
\textbf{Drawback:} Inherently suffers from aliasing, rendering it useless for highpass/bandstop filters.

\section{The Bilinear Transform (BLT)}
The BLT applies the conformal mapping:
\begin{equation}
s = \frac{2}{T_s} \frac{1 - z^{-1}}{1 + z^{-1}} = \frac{2}{T_s} \frac{z - 1}{z + 1}
\end{equation}
This maps the entire $j\Omega$ axis onto the unit circle without aliasing.
By substituting $s = j\Omega$ and $z = e^{j\omega}$, the frequency warping relationship is:
\begin{equation}
\Omega = \frac{2}{T_s} \tan\left(\frac{\omega}{2}\right)
\end{equation}

\section{Worked Example: 3rd-Order Digital Butterworth LPF}
Design a 3rd-order digital Butterworth LPF with $\omega_c = 0.3\pi$ using BLT.
\begin{enumerate}
    \item \textbf{Prewarping:} Let $T_s = 2$.
    $\Omega_c = \tan(0.15\pi) \approx 0.5095 \text{ rad/s}$.
    
    \item \textbf{Analog Prototype Design:} For $N=3$, the poles are:
    \begin{align*}
    s_1 &= 0.5095 e^{j120^\circ} = -0.2548 + j0.4412 \\
    s_2 &= 0.5095 e^{j180^\circ} = -0.5095 \\
    s_3 &= 0.5095 e^{j240^\circ} = -0.2548 - j0.4412
    \end{align*}
    The transfer function is:
    \begin{equation}
    H_a(s) = \frac{\Omega_c^3}{(s-s_1)(s-s_2)(s-s_3)} = \frac{0.1323}{s^3 + 1.0191s^2 + 0.5193s + 0.1323}
    \end{equation}
    
    \item \textbf{Bilinear Transform:} Substitute $s = \frac{z-1}{z+1}$:
    \begin{equation}
    H(z) = \frac{0.1323 (z+1)^3}{(z-1)^3 + 1.0191(z-1)^2(z+1) + 0.5193(z-1)(z+1)^2 + 0.1323(z+1)^3}
    \end{equation}
    Expanding and collecting terms yields:
    \begin{equation}
    H(z) = \frac{0.0495 (1 + 3z^{-1} + 3z^{-2} + z^{-3})}{1 - 1.1618z^{-1} + 0.6959z^{-2} - 0.1378z^{-3}}
    \end{equation}
\end{enumerate}

\section{Checkpoints}
\begin{enumerate}
    \item \textbf{Q: Why use BLT for highpass filters instead of Impulse Invariance?} \\
    \textbf{A:} Impulse Invariance causes spectral aliasing. Since highpass filters are not bandlimited, aliasing ruins the passband. BLT maps the whole axis without aliasing.
    
    \item \textbf{Q: What happens to Butterworth poles as order $N$ increases?} \\
    \textbf{A:} The poles crowd along the circle of radius $\Omega_c$ and the ones near the $j\Omega$ axis create a steeper transition band roll-off.
    
    \item \textbf{Q: Why choose $T_s = 2$ in BLT?} \\
    \textbf{A:} It simplifies calculations; $T_s$ scales the prewarped $\Omega_c$ and then is identically cancelled out during the $s \to z$ substitution, having no effect on final coefficients.
\end{enumerate}

\end{document}
